\documentclass[a4paper,12pt]{article}
\usepackage[utf8]{inputenc}
\usepackage[brazilian]{babel}
\usepackage{amsmath}
\usepackage{graphicx}
\usepackage{float}
\usepackage{hyperref}
\usepackage{amsfonts}
\usepackage{amssymb}
\usepackage{geometry}
\usepackage{listings}
\usepackage{xcolor}
% Pacote para códigos em Python
\usepackage{listings}
\usepackage[Algoritmo]{algorithm}
\usepackage{algpseudocode}
\usepackage{caption}
% \renewcommand{\lstlistingname}{Algoritmo}
% \makeatletter
% \renewcommand{\ALG@name}{Algoritmo}
% \makeatother
% Configuração do ambiente listings para Python
\lstset{
    language=Python,
    basicstyle=\ttfamily\small,
    keywordstyle=\color{blue}\bfseries,
    stringstyle=\color{red},
    commentstyle=\color{gray}\itshape,
    numbers=left,
    numberstyle=\tiny\color{gray},
    stepnumber=1,
    numbersep=5pt,
    backgroundcolor=\color{white},
    frame=single,
    tabsize=4,
    showspaces=false,
    showstringspaces=false,
    breaklines=true,
    breakatwhitespace=true,
    captionpos=b
}
\pagestyle{empty}

\geometry{a4paper, margin=0.5in}

\title{Gabarito da Prova 1 - Álgebra Linear para Computação}
\date{}
\newcommand{\avg}{\mathbf{\textbf{avg}}}
\newcommand{\std}{\mathbf{\textbf{std}}}
\newcommand{\vect}[1]{\mathbf{#1}}
\newcommand{\inner}[2]{\langle #1, #2 \rangle}
\newcommand{\norm}[1]{\left \| #1 \right \|}

\begin{document}
\maketitle

\begin{enumerate}
    \item  \begin{enumerate}
        \item Queremos encontrar $k$ tal que $\inner{\vect{u}}{\vect{v}} = 0$:
        \begin{align*}
            \inner{\vect{u}}{\vect{v}} & = 2(k+1) + 3(k-1)\\
            &= 2k + 2 + 3k - 3\\
            &= 5k - 1 = 0\\
            & \Rightarrow k = \frac{1}{5}
        \end{align*}
        Portanto, $k = \dfrac{1}{5}$.

        \item Queremos encontrar $k$ tal que $\inner{\vect{u}}{\vect{v}} = 0$:
        \begin{align*}
            \inner{\vect{u}}{\vect{v}} & = k^2 - k - 6 = 0
        \end{align*}
        Resolvendo a equação quadrática:
        \begin{align*}
            k & = \frac{1 \pm \sqrt{1 + 24}}{2}\\
            & = \frac{1 \pm 5}{2}
        \end{align*}
        Portanto, $k = 3$ ou $k = -2$.
    \end{enumerate}
    \item \begin{enumerate}
        \item A $i$--ésima entrada de $\vect{z}$ é dada por $z_i = \alpha_1 x_{1,i} + \alpha_2 x_{2,i} + \cdots +\alpha_k x_{k,i}$ onde $x_{j,i}$ representa a entrada $i$ do vetor $\vect{x}_j$. Portanto:
        \begin{align*}
            \avg(\vect{z}) & = \frac{z_1 + \cdots + z_n}{n}\\
            & = \frac{(\alpha_1 x_{1,1} + \cdots + \alpha_k x_{k,1}) + \cdots +(\alpha_1 x_{n,k} + \cdots + \alpha_k x_{k,n})}{n}\\
            & = \alpha_1 \frac{x_{1,1} + \cdots + x_{1,n}}{n} + \alpha_2 \frac{x_{2,1} + \cdots + x_{2,n}}{n} + \cdots + \alpha_k \frac{x_{k,1} + \cdots + x_{k,n}}{n}\\
            &= \alpha_1 \avg(\vect{x}_1) + \cdots + \alpha_k \avg(\vect{x}_k)
        \end{align*}
    \item Utilizando a notação $\tilde{\vect{x}_i} = \vect{x}_i - \avg(\vect{x}_i)\vect{1}$ para os vetores $\vect{x}_i$ menos as médias, temos:
    \begin{align*}
        \|\vect{z} - \avg(\vect{z})\vect{1}\|^2 & = \norm{\alpha_1\vect{x}_1 + \cdots + \alpha_k \vect{x}_k - (\alpha_1 \avg(\vect{x}_1) + \cdots + \alpha_k \avg(\vect{x}_k))\vect{1}}^2\\
        &= \norm{\alpha_1\tilde{\vect{x}_1} + \cdots + \alpha_k\tilde{\vect{x}_k}}^2\\
        &= (\alpha_1\tilde{\vect{x}_1} + \cdots + \alpha_k\tilde{\vect{x}_k})^T(\alpha_1\tilde{\vect{x}_1} + \cdots + \alpha_k\tilde{\vect{x}_k}).
    \end{align*}
    Como os vetores não são correlacionados, temos $\tilde{\vect{x}_i}^T\tilde{\vect{x}_j} = 0$ para $i \neq j$. Logo:
    \begin{align*}
        \|\vect{z} - \avg(\vect{z})\vect{1}\|^2 & =(\alpha_1\tilde{\vect{x}_1} + \cdots + \alpha_k\tilde{\vect{x}_k})^T(\alpha_1\tilde{\vect{x}_1} + \cdots + \alpha_k\tilde{\vect{x}_k}).\\
        &= \alpha_1^2\norm{\tilde{\vect{x}_1}}^2 + \cdots  \alpha_k^2\norm{\tilde{\vect{x}_k}}^2,
    \end{align*}
    portanto:
    \begin{align*}
        \std(\vect{z})^2 &= \frac{1}{n} \norm{\vect{z} - \avg(\vect{z})\vect{1}}^2\\
         &= \alpha_1^2 \frac{\norm{\tilde{\vect{x}_1}}}{n} + \cdots + \alpha_k^2 \frac{\norm{\tilde{\vect{x}_k}}}{n}\\
        & = \alpha_1^2 \std(\vect{x}_1) + \cdots +\alpha_k \std(\vect{x}_k)\\
        \Rightarrow \std(\vect{z}) & = \sqrt{\alpha_1^2 \std(\vect{x}_1) + \cdots +\alpha_k \std(\vect{x}_k)}.
    \end{align*}
    \end{enumerate}
    \item Seja $A$ a matriz cujas colunas sejam formadas pelos vetores da base $B$. Temos:
    $$
    A = \begin{bmatrix}
        a & 1 & 0 \\
        1 & a & 1 \\
        0 & 1 & a \\
    \end{bmatrix}.
    $$
    Para que $B$ Seja uma base do $\mathbb{R}^3$, temos que $\det(A) \neq 0$, ou seja:
    $$
    \det(A) = a^3 - 2a = a(a^2-2) \neq 0.
    $$
    Temos então que $\det(A) = 0$ para os valores $a = 0$, $a= \sqrt{2}$ e $a = -\sqrt{2}$, portanto $B$ é base para $a \in \mathbb{R}\setminus\{0,\sqrt{2}, -\sqrt{2}\}$.
    \item A base canônica de $\mathbb{P}_1$ é dada por $\{1, x\}$. Aplicando $T$ na base canônica, temos $T(1) = \begin{bmatrix}
        1 \\ 1
    \end{bmatrix}$ e $T(x) = \begin{bmatrix}
        0 \\ 1
    \end{bmatrix}$, portanto $[T] = \begin{bmatrix}
        1 & 0 \\ 1 & 1
    \end{bmatrix}$. Para transformação $S$, temos $S((1,0)) = \begin{bmatrix}
        1\\2
    \end{bmatrix}$ e $S((0,1)) = \begin{bmatrix}
        -2\\-1
    \end{bmatrix}$. Portanto $[S] = \begin{bmatrix}
        1 & -2\\2 & -1
    \end{bmatrix}$. Desse modo, temos $[S \circ T] = \begin{bmatrix}
        1 & 0 \\ 1 & 1
    \end{bmatrix}  \begin{bmatrix}
        1 & -2\\2 & -1
    \end{bmatrix} =  \begin{bmatrix}
        -1 & -2\\1 & -1
    \end{bmatrix}$.

    \item  Dada as funções \texttt{cria\_vetor\_canonico(i,n)}, que retorna o $i$--ésimo vetor da base canônica de $\mathbb{R}^n$, e \texttt{cria\_matriz($A_n$)}, que cria uma matriz cujas colunas são dados pelos vetores no conjunto $A_n$, temos:
    \begin{algorithm}
    \caption{Cria matriz da transformação}
    \begin{algorithmic}[1]
        \State \textbf{Entradas:} Transformação Linear $f$, dimensão de entrada $n$, dimensão do espaço de saída $m$.
        \State \textbf{Inicializa} $A_0 \gets \emptyset$
        \For{$i = 1,\cdots, n$}
            \State $e_i$ = \texttt{cria\_vetor\_canonico(i,n)}
            \State $A_{i} = A_{i-1} \cup \{f(e_i)\}$
        \EndFor
        \State $A \gets$ \texttt{cria\_matriz($A_n$)}
        \State \textbf{Retorna} $A$
    \end{algorithmic}
    \end{algorithm}

    \item \begin{enumerate}
        \item A linha 5 está calculando a projeção do vetor $v_k$ sobre os vetores $u_i$ esomando os resultados.
        \item Este é o algorítimo de ortonormalização de Gram-Schmidt. Sua finalidade é gerar uma base ortonormal a partir de um conjunto de vetores linearmente independentes. Caso o conjunto fornecido seja linearmente dependente, o vetor que é dependente dos demais é ignorado.
    \end{enumerate}
\end{enumerate}
\end{document}